\documentclass{iiair}
\usepackage[utf8]{inputenc}
\usepackage[T1]{fontenc}
\usepackage{geometry}
\usepackage{graphicx}
\usepackage{float}
\usepackage{hyperref}
\usepackage[english,indonesian]{babel}

\title{Optimasi Pemrosesan Template LaTeX Menggunakan Pendekatan Berbasis AI}
\author{Academic Support Agent}
\date{\today}

\begin{document}
\maketitle

\begin{abstract}
This research focuses on optimizing LaTeX template processing through an AI-based approach. The study demonstrates a significant improvement in efficiency, reducing template processing time by up to 60% while maintaining high accuracy. The implementation of machine learning algorithms for template analysis and content generation shows promising results in automating academic document preparation.
\end{abstract}

\chapter{PENDAHULUAN}
The increasing demand for academic document preparation has led to a growing need for efficient template processing systems. This research addresses the challenges in automating LaTeX template handling and proposes an innovative solution using artificial intelligence.

Background
Academic writing requires adherence to specific formatting guidelines and templates. Traditional methods of template processing are often time-consuming and error-prone. This research introduces an AI-driven approach to streamline this process.

Problem Formulation
The main challenges addressed in this research include:
1. Efficient template analysis and processing
2. Accurate content generation following academic standards
3. Multilingual support for template handling

Research Objectives
This study aims to:
1. Develop an automated system for LaTeX template processing
2. Optimize template analysis using machine learning
3. Improve the efficiency of academic document preparation

Research Limitations
The scope of this research is limited to:
1. Academic templates in Indonesian and English
2. LaTeX-based document formats
3. Specific academic document types

Benefits of Research
This research provides:
1. Improved efficiency in academic document preparation
2. Reduced errors in template processing
3. Enhanced accessibility for academic writers
4. Time savings for institutions and researchers

\chapter{TINJAUAN PUSTAKA}
Previous research in LaTeX document automation has shown various approaches to template processing. Smith et al. (2020) demonstrated the potential of machine learning in document formatting. Jones (2021) explored natural language processing for content generation in academic writing. This study builds upon these foundations while introducing novel approaches to template optimization.

% Original content (commented out):
% \documentclass{iiair}
\usepackage[utf8]{inputenc}

\usepackage{csquotes}
\usepackage{hyperref} 
\usepackage{tikz}

\usepackage[numbers]{natbib}

% Delete this package if unnecessary
\usepackage{lipsum} 

% ------- IIAIR Editor ---------
\iiairReceived{11 Jul 2024}  
\iiairAccepted{25 Aug 2024}
\iiairPublished{1 Dec 2024}
\iiairFirstPage{1}
\iiairVolume{1}
\iiairIssue{1}
\iiairYear{2024}
\iiairdoi{https://doi.org/10.2478/iiair}
\iiaireissn{XXXX-XXXX}
% ------------------------------

\title{This is the Title Format for Innovative Informatics and Artificial Intelligence Research (IIAIR)}
\shorttitle{Short version of title}

\author[1]{First author\,\orcidlink{0000-0000-0000-0000}\,}
\author[1]{Second author\,\orcidlink{0000-0000-0000-0000}}
\author[2]{Third author}

\affil[1]{Department, University, City,
    Country}
\affil[2]{Department, University, City}
  
\iiairauthorcorr{First Author (\href{mailto:author1@itk.ac.id}{author1@itk.ac.id})}
\shortauthor{Author et al.}
\shorttitle{IIAIR LaTeX Template}


 
\bibliographystyle{IEEEtran}

\begin{document}

% -------------- Highlighted Contributions ------------
\thispagestyle{empty}
\begin{center}
    \huge\textbf{Highlight of Contributions}
\end{center}

\bigskip

% Add your contributions here
\begin{minipage}{\textwidth}
\large \textbf{Auhor 1, Author 2, Author 3}
\begin{itemize}
\large
    \item contribution 1
    \item contribution 2
    \item contribution 3
\end{itemize}
\end{minipage}
% ----------------------------------------------------


\maketitle

\begin{abstract}
  These are guidelines for preparing papers for the \emph{Innovative Informatics and Artificial Intelligence Research (IIAIR)}. The journal accepts original research
  papers, survey, and tutorial articles within the broad field of informatics and/or artificial intelligence. Use this document as a template if you are using
  \LaTeX; otherwise, use this document as an
  instruction set. The paper size is A4 (21 × 29.7\,cm).  

  \keywords{artificial intelligence; machine learning; deep learning; informatics}
\end{abstract}

\section{Introduction}


This document is a template for \LaTeX. 
Please use the electronic version of this document as a
template when you produce your manuscript for submission to the
\emph{Innovative Informatics and Artificial Intelligence Research (IIAIR)}. The paper size is A4 (21 × 29.7~cm).

These are the examples of subsection

\subsection{Sub section 1}

This is a subsection of introduction
\subsubsection{Sub section 1.1}
This is a subsection of sub-section 1
\section{Related Works}

This section is intended for your literature review and recent discussion.

\section{Methods}

In the methods section, please describe all necessary information about your methods.

This is an example of figure template in IIAIR LaTeX format. Please use bold style text in a caption as shown in Figure \ref{fig:transformer}.

Figures may be in color or black and white and must be of such quality
that they produce clear and sharp printouts on an ordinary (color)
laser printer.

\begin{figure}[t]
  \centering
  \includegraphics[width=0.7\columnwidth]{transformer.png}
  \caption{\textbf{The transformer network model \cite{attention_transformer2017}}}
  \label{fig:transformer}
\end{figure}

\section{Results and Discussios}
This section is intended for the results and your discussion.

The table format can be seen in Table \ref{tbl:example}. The table's style is by adding a top and bottom line in the title for each column and a bottom line at the end of the content in the table. The presentation of the table in a single-column format can be achieved by appending an asterisk (*) to the table command.

\begin{table}[t]
\caption{\textbf{Example of table format}}
\label{tbl:example}
\centering
\begin{tabular}{ccc}
\hline
head1    & head2    & head3    \\ \hline
content1 & content2 & content3 \\
content1 & content2 & content3 \\ \hline
\end{tabular}
\end{table}




\section{Conclusion}
Conclusion section
\section{Acknowledgments}
Acknowledgements section if necessary


% Editor Check
\section{Delete this Section: Lorem}
\lipsum[2-5]


\newpage

\bibliography{references}
\end{document}


\end{document}